\Needspace{8\baselineskip}
\Article{289}
Aus der allgemeinen Theorie in Art.\ 282 folgt für die eigentlichen
Darstellungen binärer Formen durch die ternäre Form $xx+yy+zz=f$ Folgendes.

I.\ Eine binäre Form $\varphi$ kann nur dann durch $f$ eigentlich
dargestellt werden, wenn sie eine primitive positive Form ist und
$-1$ (also die Determinante der Form $f$) ihre charakteristische
Zahl ist. Daher gibt es bei positiver Determinante und ebenso bei
negativer Determinante $-M$, wenn $M$ entweder durch $4$
teilbar oder von der Form $8n+7$ ist, keine durch $f$ eigentlich
darstellbaren binären Formen.

II.\ Ist dagegen $\varphi=(p,q,r)$ eine primitive positive Form mit der
Determinante $-M$ und $-1$ eine charakteristische Zahl der
Form $\varphi$, also auch der entgegengesetzten Form $(p,-q,r)$, so gibt
es eigentliche Darstellungen der Form $\varphi$ durch $f$, die zu
jedem gegebenen Wert des Ausdrucks $\surd{}-(p,-q,r)$ gehören. Denn alle Koeffizienten
der ternären Form $g$ mit der Determinante $-1$ (Art.\ 283)
sind notwendig ganzzahlig, und $g$ ist eine definite Form, also
sicher zu $f$ äquivalent (Art.\ 285, I).

III.\ Nach Art.\ 283, III ist die Anzahl aller Darstellungen, die
zu demselben Wert des Ausdrucks $\surd{}-(p,-q,r)$ gehören, in allen Fällen
außer $M=1$ und $M=2$ ebenso groß wie die Anzahl der Transformationen
der Form $f$ in $g$, also nach Art.\ 285 $=48$. Daraus
erkennt man auch: Hat man eine zu einem gegebenen Wert gehörende
Darstellung, so erhält man die übrigen $47$, indem man die Werte
von $x$, $y$, $z$
\SourceLeaf{346}
auf alle möglichen Arten miteinander vertauscht und ihre Vorzeichen
umkehrt. Alle $48$ Darstellungen ergeben daher \textit{eine einzige}
Zerlegung der Form $\varphi$ in drei Quadrate, wenn man nur die Quadrate
selbst betrachtet, ohne ihre Reihenfolge oder die Vorzeichen ihrer
Wurzeln zu berücksichtigen.

IV.\ Die Anzahl aller verschiedenen ungeraden Primzahlen, die $M$
teilen, sei $=\mu$. Aus Art.\ 233 folgt leicht, dass die Anzahl
aller verschiedenen Werte des Ausdrucks $\surd{}-(p,\allowbreak -q,\allowbreak r)(\mathrm{mod.}\ M)$ $=2^\mu$ ist.
Nach Art.\ 283 braucht man davon nur die Hälfte zu betrachten
(wenn $M>2$). Die Anzahl aller eigentlichen Darstellungen der
Form $\varphi$ durch $f$ ist daher $=48\,\allowbreak  .\,\allowbreak  2^{\mu-1}=3\,\allowbreak  .\,\allowbreak  2^{\mu+3}$, die Anzahl der verschiedenen
Zerlegungen in drei Quadrate aber $=2^{\mu-1}$.

\textit{Beispiel.} Sei $\varphi=19tt+6tu+41uu$, also $M=770$. Hier muss man die folgenden
vier Werte des Ausdrucks $\surd{}-(19,\allowbreak -3,\allowbreak 41)(\mathrm{mod.}\ 770)$ betrachten (Art.\ 283): $(39,237)$,
$(171,-27)$, $(269,-83)$, $(291,-127)$. Um die zum Wert $(39,237)$ gehörenden Darstellungen
zu finden, bestimmt man zuerst die ternäre Form $\left(\begin{smallmatrix}19,&41,&2\\3,&6,&3\end{smallmatrix}\right)=g$. Nach den
Regeln aus Art.\ 272, 275 geht $f$ durch die Substitution $\left\{\begin{smallmatrix}1,&-6,&-0\\-3,&-2,&-1\\-3,&-1,&-1\end{smallmatrix}\right\}$
in sie über. Daraus erhält man die folgende Darstellung der Form
$\varphi$ durch $f$:
\[
x=t-6u,\quad y=-3t-2u.\quad z=-3t-u
\]
Die übrigen $47$ zum selben Wert gehörenden Darstellungen, die
durch Vertauschen dieser Werte und Umkehren ihrer Vorzeichen entstehen,
schreiben wir der Kürze wegen nicht auf. Alle $48$ Darstellungen
ergeben dieselbe Zerlegung der Form $\varphi$ in die drei Quadrate
\[
tt-12tu+36uu,\quad 9tt+12tu+4uu,\quad 9tt+6tu+uu
\]

Auf genau dieselbe Weise liefert der Wert $(171,-27)$ die Zerlegung in
die Quadrate $(3t+5u)^2$, $(3t-4u)^2$, $tt$; der Wert $(269,-83)$ liefert $(t+6u)^2+(3t+u)^2+(3t-2u)^2$
und schließlich der Wert $(291,-127)$ die Zerlegung $(t+3u)^2+(3t+4u)^2+(3t-4u)^2$. Jede dieser
Zerlegungen entspricht $48$ Darstellungen. --- Außer diesen $192$
Darstellungen oder vier Zerlegungen gibt es keine weiteren, da $770$
durch kein Quadrat teilbar ist und somit keine uneigentlichen Darstellungen
existieren können.

\SourceLeaf{347}
\Needspace{8\baselineskip}
\Article{290}
Die Formen mit den Determinanten $-1$ und $-2$, für die gewisse
Ausnahmen galten, besprechen wir kurz gesondert. Zunächst die allgemeine
Bemerkung: Sind $\varphi$, $\varphi'$ beliebige äquivalente binäre Formen
und $(\Theta)$ eine gegebene Transformation der ersten in die zweite,
so ergibt die Verbindung einer beliebigen Darstellung der Form
$\varphi$ durch eine ternäre Form $f$ mit der Substitution $(\Theta)$
eine Darstellung der Form $\varphi'$ durch $f$. Dabei entstehen aus
eigentlichen Darstellungen von $\varphi$ eigentliche Darstellungen
von $\varphi'$, aus verschiedenen verschiedene und schließlich aus
allen sämtliche. All dies lässt sich sehr leicht durch Rechnung
bestätigen. Die eine der Formen $\varphi$, $\varphi'$ lässt sich also auf
genauso viele Arten durch $f$ darstellen wie die andere.

I.\ Sei zuerst $\varphi=tt+uu$ und $\varphi'$ eine beliebige andere positive
binäre Form mit der Determinante $-1$, zu der $\varphi$ daher
äquivalent ist. Es gehe $\varphi$ in $\varphi'$ durch die Substitution
$t=\alpha t'+\beta u'$, $u=\gamma t'+\delta u'$ über. Die Form $\varphi$ wird durch die ternäre Form $f=zz+yy+zz$
dargestellt, wenn man $x=t$, $y=u$, $z=0$ setzt. Durch Vertauschen
von $x,y,z$ entstehen daraus \textit{sechs} Darstellungen und aus
jeder durch Ändern der Vorzeichen von $t$, $u$ wiederum
vier. So erhält man insgesamt $24$ verschiedene Darstellungen,
denen eine einzige Zerlegung in drei Quadrate entspricht; dass es
keine weiteren Darstellungen geben kann, sieht man leicht. Daher
lässt sich auch die Form $\varphi'$ nur auf eine einzige Weise in
drei Quadrate zerlegen, nämlich in $(\alpha t'+\beta u')^2$, $(\gamma t'+\delta u')^2$ und $0$.
Diese Zerlegung entspricht $24$ Darstellungen.

II.\ Sei $\varphi=tt+2uu$ und $\varphi'$ eine beliebige andere positive binäre
Form mit der Determinante $-2$, in die $\varphi$ durch die
Substitution $t=\alpha t'+\beta u'$, $u=\gamma t'+\delta u'$ übergeht. Dann folgt wie im vorigen
Fall, dass sich $\varphi$ und somit auch $\varphi'$ nur auf eine einzige
Weise in drei Quadrate zerlegen lässt: $\varphi$ in $tt+uu+uu$ und
$\varphi'$ in $(\alpha t'+\beta u')^2+(\gamma t'+\delta u')^2+(\gamma t'+\delta u')^2$. Dass eine solche Zerlegung $24$ Darstellungen
entspricht, sieht man leicht.

Daraus folgt, dass die binären Formen mit den Determinanten $-1$
und $-2$ hinsichtlich der Anzahl ihrer Darstellungen durch die
ternäre Form $xx+yy+zz$ mit den anderen binären Formen völlig übereinstimmen.
Denn in beiden Fällen ist $\mu=0$, und die Formel aus IV des vorigen
Artikels ergibt tatsächlich $24$ Darstellungen. Der Grund dafür
ist, dass sich die beiden Ausnahmen, denen solche Formen unterlagen,
gegenseitig ausgleichen.

Die Anwendung der in Art.\ 284 entwickelten allgemeinen Theorie
der uneigentlichen Darstellungen auf die Form $xx+yy+zz$ lassen wir
der Kürze wegen weg.

\SourceLeaf{348}
\Needspace{8\baselineskip}
\Article{291}
Die Aufgabe, alle eigentlichen Darstellungen einer gegebenen positiven
\textit{Zahl} $M$ durch die Form $xx+yy+zz$ zu finden, wird nach
Art.\ 281 zunächst auf die Suche nach den eigentlichen Darstellungen
der Zahl $-M$ durch die Form $-xx-yy-zz=f$ zurückgeführt. Diese erhält
man nach den Regeln aus Art.\ 280 wie folgt.

I.\ Man bestimme alle Klassen binärer Formen mit der Determinante
$-M$, deren Formen sich durch $XX+YY+ZZ=F$ (die ternäre Form, deren
Adjunkte $f$ ist) eigentlich darstellen lassen. Ist $M\equiv0,\allowbreak \ 4\text{ oder }7\ (\mathrm{mod.}\ 8)$,
so gibt es nach Art.\ 288 keine solchen Klassen, und $M$ lässt
sich daher nicht in drei Quadrate ohne gemeinsamen Teiler
zerlegen\HistoricalNote{*)}{Diese Unmöglichkeit folgt auch daraus,
dass die Summe dreier ungerader Quadrate notwendig $\equiv3\ (\mathrm{mod.}\ 8)$ ist;
die Summe zweier ungerader und eines geraden Quadrats ist entweder
$\equiv2$ oder $\equiv6$; die Summe eines ungeraden und zweier gerader
Quadrate entweder $\equiv1$ oder $\equiv5$; schließlich die Summe
dreier gerader Quadrate entweder $\equiv0$ oder $\equiv4$. Im letzten
Fall ist die Darstellung aber offenbar uneigentlich.}. Ist dagegen
$M\equiv1,\ 2,\ 5\text{ oder }6$, so gibt es ein positives eigentlich primitives Geschlecht,
und bei $M\equiv3$ ein uneigentlich primitives, das alle diese Klassen
umfasst. Ihre Anzahl bezeichnen wir mit $k$.

II.\ Nun wähle man aus diesen Klassen beliebig $k$ Formen,
aus jeder eine, und bezeichne sie mit $\varphi$, $\varphi'$, $\varphi''$ usw.
Man bestimme sämtliche eigentlichen Darstellungen aller dieser Formen
durch $F$. Ihre Anzahl ist daher $3\, .\, 2^{\mu+3}k=K$, wobei $\mu$ die
Anzahl der (ungeraden) Primfaktoren von $M$ bezeichnet. Schließlich
leite man aus jeder solchen Darstellung, etwa
\[
X=mt+nu,\quad Y=m't+n'u,\quad Z=m''t+n''u
\]
die folgende Darstellung von $M$ durch $xx+yy+zz$ ab:
\[
x=m'n''-m''n',\quad y=m''n-mn'',\quad z=mn'-m'n
\]
Die Gesamtheit dieser $K$ Darstellungen, die wir mit $\Omega$
bezeichnen, enthält notwendig sämtliche Darstellungen von $M$.

III.\ Es bleibt also nur zu untersuchen, ob in $\Omega$ \textit{identische}
Darstellungen auftreten können. Nach Art.\ 280, III sind die Darstellungen
in $\Omega$, die aus verschiedenen Formen, \textit{zum Beispiel}
aus $\varphi$ und $\varphi'$, hergeleitet wurden, notwendig
\SourceLeaf{349}
verschieden. Es bleibt daher nur die Frage, ob verschiedene Darstellungen
derselben Form, \textit{zum Beispiel} von $\varphi$, durch $F$
identische Darstellungen der Zahl $M$ durch $xx+yy+zz$ ergeben
können. Nun ist sofort klar: Kommt unter den Darstellungen von $\varphi$
die folgende vor,
\[
X=mt+nu,\quad Y=m't+n'u,\quad Z=m''t+n''u\quad\ldots\ (r)
\]
so auch diese:
\[
X=-mt-nu,\quad Y=-m't-n'u,\quad Z=-m''t-n''u\quad\ldots\ (r')
\]
Aus beiden folgt dieselbe Darstellung von $M$, die wir mit
$(R)$ bezeichnen. Wir untersuchen also, ob sich dieselbe $(R)$
noch aus anderen Darstellungen der Form $\varphi$ ergeben kann.
Setzt man in Art.\ 280, III $\chi=\varphi$, so folgt leicht: Werden alle
eigentlichen Transformationen der Form $\varphi$ in sich selbst durch
\[
t=\alpha t+\beta u,\quad u=\gamma t+\delta u
\]
gegeben, so werden alle Darstellungen der Form $\varphi$, aus denen
$R$ folgt, durch
\[
\begin{aligned}
x&=(\alpha m+\gamma n)t+(\beta m+\delta n)u\\
y&=(\alpha m'+\gamma n')t+(\beta m'+\delta n')u\\
z&=(\alpha m''+\gamma n'')t+(\beta m''+\delta n'')u
\end{aligned}
\]
ausgedrückt. Aus der in Art.\ 179 entwickelten Theorie der Transformationen
binärer Formen mit negativer Determinante folgt aber, dass es in
allen Fällen außer $M=1$ und $M=3$ nur zwei eigentliche Transformationen
der Form $\varphi$ in sich selbst gibt, nämlich $\alpha,\allowbreak \beta,\allowbreak \gamma,\allowbreak \delta=1,\allowbreak 0,\allowbreak 0,\allowbreak 1$ und $=-1,0,0,-1$.
(Da $\varphi$ eine primitive Form ist, ist das in Art.\ 179 mit
$m$ Bezeichnete entweder $1$ oder $2$. Abgesehen von
den Ausnahmefällen gilt dort also sicher $(1)$.) Somit kann
$(R)$ nur aus $r$, $r'$ entstehen. Jede eigentliche
Darstellung der Zahl $M$ kommt daher in $\Omega$ zweimal und
nicht öfter vor, und die Anzahl aller verschiedenen eigentlichen
Darstellungen von $M$ ist $\tfrac{1}{2}K=3\,\allowbreak  .\,\allowbreak  2^{\mu+2}k$.

Für die Ausnahmefälle ist die Anzahl der eigentlichen Transformationen
der Form $\varphi$ in sich selbst nach Art.\ 179 gleich $4$
für $M=1$ und $6$ für $M=3$. Tatsächlich bestätigt man
leicht, dass die Anzahlen der eigentlichen Darstellungen der Zahlen
$1$, $3$ gleich $\tfrac{1}{4}K$ bzw.\ $\tfrac{1}{6}K$ sind. Jede dieser
Zahlen lässt sich nämlich nur auf eine einzige Weise in drei
\SourceLeaf{350}
Quadrate zerlegen: $1$ in $1+0+0$ und $3$ in $1+1+1$.
Die Zerlegung von $1$ liefert sechs, die von $3$ acht
verschiedene Darstellungen. Dabei ist $K$ $=24$
für $M=1$ (wobei $\mu=0$, $k=1$) und $=48$ für
$M=3$ (wobei $\mu=1$, $k=1$).

Außerdem bemerken wir: Bezeichnet $h$ die Anzahl der Klassen
im Hauptgeschlecht, der nach Art.\ 252 die Anzahl der Klassen in
jedem anderen eigentlich primitiven Geschlecht gleich ist, so ist
$k=h$ für $M\equiv1,\allowbreak \ 2,\allowbreak \ 5\text{ oder }6\ (\mathrm{mod.}\ 8)$, aber $k=\tfrac{1}{3}h$ für $M\equiv3\ (\mathrm{mod.}\ 8)$, mit der einzigen
Ausnahme $M=3$, wo $k=h=1$ ist. Für Zahlen der Form $8n+3$
ist die Anzahl der Darstellungen also \textit{allgemein} $=2^{\mu+2}h$,
da sich bei der Zahl $3$ die beiden Ausnahmen gegenseitig ausgleichen.
